\section{Watermarking and removal
experiments}\label{record-docs-watermarking-md-watermarking-and-removal-experiments}

The current text protocol is
\href{records/docs/instruct_watermarking.tex}{TextSeal 2 and
DistSeal-inspired instruct Qwen}. Earlier TextSeal 1 and Base-model text
outputs and metrics have been deleted. The following describes preserved
DistSeal image experiments.

The new \href{records/docs/closed_image_watermarking.tex}{closed-source
image protocol} instead tests the released NeurIPS 2025 NFPA attack
using only image outputs, with an independent public regeneration model
and separately held owner detectors.

\subsection{Image model
removal}\label{record-docs-watermarking-md-image-model-removal}

Released RAR-XL and UViT-H/DC-AE generators are evaluated with their
original unmodified decoders and extractors. Clean-checkpoint
interpolation at fractions 0.75, 0.5, 0.25 and full replacement provides
reference-weight controls. Clean LoRA and full-parameter training at 500
and 2,500 steps test removal without clean reference weights. Owner keys
and scores do not enter training. Task schemas, separate processes and
artifact hashes enforce input contracts; they are not an
operating-system sandbox against malicious workers.

Image training uses 8,000 clean ImageNet training images from 800
classes and 800 separate images for validation. Classes 800 through 999
are excluded from attack fitting; owner evaluation reports these
separately. Labels follow the original fixed synset order. The ordinary
generator backbone is the predefined training target; the decoder is
unchanged without consulting watermark location.

\subsection{Relationship to
DistSeal}\label{record-docs-watermarking-md-relationship-to-distseal}

\href{https://proceedings.mlr.press/v306/rebuffi26a.html}{DistSeal}
trains a latent embedder/extractor and then trains a generator on its
watermarked latents, or trains a decoder to reconstruct watermarked
images. Appendix G of the ICML 2026 version demonstrates forgetting
during clean LoRA fine-tuning, using 50 images of one ImageNet class.
The earlier single-class pilot uses that setting, with 2,500 LoRA
removal steps at learning rate 1e-4. The later primary protocol instead
fits on 8,000 images from 800 classes and uses LoRA learning rate 1e-4
or full-update learning rate 1e-5. Both retain 500-step and 2,500-step
outputs. Evaluation separates fitting and unseen classes.

\subsection{Removal literature and
interpretation}\label{record-docs-watermarking-md-removal-literature-and-interpretation}

\href{https://arxiv.org/abs/2502.10525}{Towards Watermarking of
Open-Source LLMs} systematically evaluates merging, clean fine-tuning,
pruning and quantization. Its results motivate testing weight changes
while measuring model utility, and evaluating the domain targeted by
fine-tuning. Our bounded experiments use linear interpolation, rather
than reproducing that paper's SLERP implementation.

Meta's \href{https://arxiv.org/abs/2303.15435}{Stable Signature} embeds
watermarks in the latent decoder and studies reconstruction-based
purification and collusion.
\href{https://arxiv.org/abs/2405.07145}{Stable Signature is Unstable}
further studies decoder removal with and without encoder access,
combining reconstruction, perceptual and adversarial losses. Those
decoder attacks are relevant context; the jobs here target DistSeal's
watermarked generators using their native clean training objectives.
They do not implement that decoder attack as if it were a DistSeal
reproduction.

Replacing a model with its exact original trivially removes changes
introduced by watermark training. Intermediate interpolation points
measure whether removal also succeeds while retaining some modified
weights. Withholding an original copy does not establish security: clean
fine-tuning may still erase a watermark. Published image checkpoint
messages are public, so these runs test removal rather than secret-key
security.

\subsection{Image
measurements}\label{record-docs-watermarking-md-image-measurements}

The historical pilot uses 50 ImageNet training images from class zero;
the primary removal run uses the 8,000-image fitting protocol above.
Each generated evaluation uses 512 images, including 128 from class
zero; calibration uses 256 different base-model noise seeds. Metrics
include bit accuracy, detection, false positives, confidence intervals,
AUROC and CLIP class alignment. JPEG 50 and a center crop retaining half
the area are evaluated with thresholds frozen from correspondingly
transformed calibration images. These sample counts cannot validate
extremely small false-positive rates or reproduce the paper's
50,000-image FID results.

Preserved mixed-run manifests and source snapshots retain their original
hashes for image provenance. Removed text branches must not be
interpreted as current results. Use the new text driver for new text
experiments.

For the complete public tables, method coverage and eight-pair examples,
use the \href{../docs/watermarking/README.md}{portable watermark
report}. It distinguishes these protocols and includes all recorded
JPEG/crop cross-conditions.
